% !TEX root = ../main.tex
\chapter{RCWA 在周期开放结构与阈值分析中的应用}
\label{ch:rcwa}

\section*{学习目标}
\begin{itemize}
  \item 理解为什么在 \cref{ch:gamma} 建立能带概念、并在 \cref{ch:pwe} 用 PWEM 算出候选模目录之后，还必须引入 RCWA。
  \item 掌握 RCWA 的散射矩阵视角，明白它为什么特别适合处理层状周期开放结构。
  \item 学会区分“周期单胞被动共振结论”“光学阈值增益估计”和“有限器件级结论”这三个层次。
\end{itemize}

\section*{先修知识}
建议已掌握 \cref{ch:pwe,ch:cwt-min,ch:threshold} 中关于候选模筛选、耦合波图像和阈值条件的内容。

\section*{本章路线图}
\ChapterModelLevel{L1→L2}
\cref{ch:gamma} 已经说明怎样阅读能带图，\cref{ch:pwe} 又把它落实成了 PWEM 计算，于是到这里读者已经知道“候选模在哪里、怎样随几何变化、主导平面波分量大致是什么”。但这仍然不是面发射激光器的完整光学问题，因为 PWEM 还没有真正处理开放边界，也不能直接回答法向辐射如何离开器件。本章要做的事，就是把问题推进到“周期开放结构的散射与共振”层级。我们先说明为什么 RCWA 特别适合 PCSEL，再从分层和傅里叶展开建立散射矩阵框架，随后解释共振极点、被动谱和光学阈值估计之间的关系，最后明确它的边界：RCWA 非常强，但它默认描述的仍是周期单胞，不是完整有限器件。

\section{为什么在 PWEM 之后自然轮到 RCWA}
PWEM 解决的是被动本征问题，能告诉我们 Gamma 点附近有哪些候选模、主要由哪些平面波分量组成。但 PCSEL 的关键物理之一——\textbf{模式如何向自由空间辐射}——在 PWEM 的封闭本征框架里根本没有出现。

因此，接下来必须问的不再是”有哪些本征模”，而是”入射波和出射波如何经过层状周期结构发生散射”。这正是 RCWA 的天然视角：它把周期结构看成散射系统，直接处理开放边界、入射通道与辐射耦合。对 PCSEL 来说，这一视角能直接回答三类问题：某个候选模是亮模还是暗模；某种几何扰动主要增强的是上方辐射还是下方辐射；以及在给定小信号增益下，开放极点何时靠近实轴。

\begin{IntuitionBox}
PWEM 回答”无限周期结构里有哪些候选模家族”；RCWA 则开始回答”这些模式怎样与外界交换能量”。走到这一步，”面发射”才真正进入模型主干。
\end{IntuitionBox}
\paragraph{与第 12 章的对比} PWEM 擅长快速扫描晶格参数并找到候选模频段，但它默认封闭边界，无法区分向上辐射和向下辐射的贡献。RCWA 则以开放散射问题为出发点，可以直接计算$S_{11}$（反射）和$S_{21}$（透射），从而给出上下出光比例。这就是为什么我们需要在学习 PWEM 之后再学习 RCWA。
\section{RCWA 的基本思路：平面内做傅里叶，纵向做分层}
RCWA 的名称里虽然有“耦合波”，但它和前面 CWT 的工作方式并不相同。CWT 的核心是主动保留少数主导波分量，用低维耦合模型换取解释性；RCWA 则更接近一种严格的半解析数值方法，它在每一层中对平面内周期方向做傅里叶展开，而在纵向把结构切成若干层逐层求解。

设器件在 $x$ 和 $y$ 方向周期，在 $z$ 方向分层。对每一层 $\ell$，把介电常数与场都在平面内写成傅里叶级数，就得到一个有限维模态系统。该层中的场可写成“向上”和“向下”传播的模叠加：
\begin{equation}
\bm{f}_\ell(z)
=
\mathbf{W}_\ell^+\bm{a}_\ell^+\ee^{+\ii\mathbf{\beta}_\ell z}
+
\mathbf{W}_\ell^-\bm{a}_\ell^-\ee^{-\ii\mathbf{\beta}_\ell z},
\label{eq:rcwa_layer_field}
\end{equation}
其中 $\bm{f}_\ell$ 是该层选取的场变量向量，$\mathbf{W}_\ell^{\pm}$ 是本征模基底矩阵，$\mathbf{\beta}_\ell$ 给出纵向传播常数，$\bm{a}_\ell^{\pm}$ 分别表示上、下行模振幅。

这一步的物理意义非常明确：在每一层里，平面内周期导致的多谐波耦合已经被吸收进层本征模中；而真正的结构复杂性，被转化成"层与层之间如何通过边界条件连接"。于是，复杂三维问题被拆成两部分：层内求模，层间拼接。\cref{fig:ch13_rcwa_discretization}直观展示了如何将连续的倾斜光栅剖面离散化为多层阶梯近似结构，这是 RCWA 处理任意形状周期结构的核心技巧。

% !TEX root = ../main.tex
% Figure: RCWA unit cell discretization into layers
% Chapter: ch13 - Rigorous Coupled-Wave Analysis
% Description: Staircase approximation of arbitrary profiles

\begin{figure}[htbp]
  \centering
  \resizebox{\textwidth}{!}{%
  \begin{tikzpicture}[scale=1.15]
    % Original continuous profile (slanted grating)
    \begin{scope}[shift={(-4,0)}]
      % Substrate
      \fill[gray!30] (-2,-1) rectangle (2,0);
      \draw[thick] (-2,0) -- (2,0);
      \node[right] at (2.1,-0.5) {衬底$n_s$};
      
      % Slanted grating teeth
      \foreach \x in {-1.5,-0.75,...,1.5} {
        \fill[blue!20] (\x,0) -- (\x+0.3,0) -- (\x+0.5,1.5) -- (\x+0.2,1.5) -- cycle;
        \draw[thick, blue] (\x,0) -- (\x+0.3,0) -- (\x+0.5,1.5) -- (\x+0.2,1.5) -- cycle;
      }
      
      % Superstrate (air)
      \fill[blue!5] (-2,1.5) rectangle (2,2.5);
      \draw[thick] (-2,1.5) -- (2,1.5);
      \node[right] at (2.1,2) {超strate$n_0$};
      
      % Incident plane wave
      \draw[->, very thick, red] (-2.5,1) -- (-1.8,1.3);
      \node[left, red] at (-2.5,1) {入射光};
      
      % Period indication
      \draw[<->, thick] (-1.5,-0.3) -- node[below] {$\Lambda$} (-0.75,-0.3);
      
      % Height indication
      \draw[<->, thick] (2.3,0) -- node[right] {$h$} (2.3,1.5);
      
      % Tilt angle
      \draw[->] (-1.5,0) arc (90:75:0.8);
      \node[right] at (-1.3,0.6) {\scriptsize $\theta_t$};
      
      \node[above, font=\bfseries] at (0,2.8) {(a) 真实倾斜光栅};
      
      % Fill factor label
      \node[align=center, font=\small] at (0,-0.8) {占空比$f=w/\Lambda$\\锥角$\theta_t$};
    \end{scope}
    
    % Discretized staircase approximation
    \begin{scope}[shift={(2,0)}]
      % Substrate
      \fill[gray!30] (-2,-1) rectangle (2,0);
      \draw[thick] (-2,0) -- (2,0);
      \node[right] at (2.1,-0.5) {衬底$n_s$};
      
      % Staircase layers (10 slices)
      \foreach \y [count=\i] in {0.15,0.3,...,1.5} {
        \pgfmathsetmacro{\width}{1.5 - \i*0.1}
        \fill[green!\i0!blue!20] (-1.5,\y-0.15) rectangle ++(\width,0.15);
        \draw[dashed, gray] (-2,\y) -- (2,\y);
      }
      
      % Layer numbering
      \node[left, font=\small] at (-2.1,0.75) {$\ell=1$};
      \node[left, font=\small] at (-2.1,1.5) {$\ell=N_L$};
      
      % Effective index in each layer
      \node[right, font=\footnotesize] at (1.2,0.5) {$n_{\mathrm{eff}}^{(1)}$};
      \node[right, font=\footnotesize] at (1.2,1.0) {$n_{\mathrm{eff}}^{(2)}$};
      \node[right, font=\footnotesize] at (1.2,1.4) {$n_{\mathrm{eff}}^{(N_L)}$};
      
      % Superstrate
      \fill[blue!5] (-2,1.5) rectangle (2,2.5);
      \draw[thick] (-2,1.5) -- (2,1.5);
      \node[right] at (2.1,2) {超strate$n_0$};
      
      \node[above, font=\bfseries] at (0,2.8) {(b) 阶梯近似分层};
      
      % Number of layers
      \draw[<->, thick] (2.3,0) -- node[right] {$N_L=10$} (2.3,1.5);
    \end{scope}
    
    % Connecting arrow showing approximation concept
    \draw[->, ultra thick, orange, bend left=15] (-1.5,2.5) to[out=20,in=160] 
      node[above, align=center, font=\small] {Li 因子化 + 傅里叶谐波展开\\\\tiny 每层内介电常数$\varepsilon(x,y)$分段恒定} (0,2.5);
    
    % Bottom: Fourier series representation
    \node[draw, align=center, font=\small] at (0,-1.5) {
      第$\ell$层介电函数的傅里叶展开:\\
      $\varepsilon^{(\ell)}(x) = \sum_{m=-\infty}^{\infty}\hat{\varepsilon}_m^{(\ell)}e^{i m G x}$,\\
      其中$G=2\pi/\Lambda$,截断至$m=\pm N_h$个谐波\\[4pt]
      S-matrix 算法保证数值稳定性
    };
    
    % Convergence note
    \node[align=left, font=\footnotesize, anchor=south west] at (-3.5,-2.5) {
      \textbf{收敛准则}:\\
      $\bullet$ 谐波数$N_h\geq 21$保证$<1\%$误差\\
      $\bullet$ 分层数$N_L$:陡峭侧壁需更细划分\\
      $\bullet$ Li 规则处理不连续面避免 Gibbs 震荡
    };
  \end{tikzpicture}%
  }
  \caption{RCWA 中常用的阶梯离散化策略。(a) 上图示意具有倾斜侧壁的真实周期结构；(b) 下图示意将其沿纵向切分为 $N_L$ 个矩形薄层，并在每层内用有限谐波数 $N_h$ 展开介电函数。配合界面连续条件和散射矩阵递推，可得到整体的反射、透射与模式耦合信息。}
  \label{fig:ch13_rcwa_discretization}
\end{figure}


\section{为什么散射矩阵是 RCWA 的真正核心}
如果只是逐层传递振幅，最直接的想法是用传输矩阵（transfer matrix）。但对高对比度、强倏逝波或厚层结构，这样做在数值上常常不稳定，因为指数增长和指数衰减模式会在矩阵乘法中导致严重病态。RCWA 在工程上真正稳健的做法，是使用散射矩阵（scattering matrix）递推。

把“不稳定”说得更定量一点：若某层的纵向本征常数写成 $k_z=\ii\kappa$，则该层传播因子会同时含有
\[
\ee^{+\ii k_z d}=\ee^{-\kappa d},
\qquad
\ee^{-\ii k_z d}=\ee^{+\kappa d}.
\]
在简单传输矩阵连乘中，指数衰减项和指数增长项会被强行放进同一矩阵乘法链，结果条件数往往随层数和倏逝强度指数恶化。散射矩阵递推的价值就在于：它把入射、反射和透射通道分开追踪，避免把“物理上应当衰减的分量”和“数值上会炸掉的增长因子”直接混在同一个乘法序列里。

把整个结构视为一个多端口散射体，就有
\begin{equation}
\begin{bmatrix}
\bm{b}_{\mathrm{up}}\\
\bm{b}_{\mathrm{down}}
\end{bmatrix}
=
\mathbf{S}(\omega,\kvec_\parallel)
\begin{bmatrix}
\bm{a}_{\mathrm{up}}\\
\bm{a}_{\mathrm{down}}
\end{bmatrix},
\label{eq:rcwa_smatrix}
\end{equation}
其中 $\bm{a}_{\mathrm{up}}$ 与 $\bm{a}_{\mathrm{down}}$ 表示从上、下方入射的通道振幅，$\bm{b}_{\mathrm{up}}$ 与 $\bm{b}_{\mathrm{down}}$ 则是对应的出射振幅。$\mathbf{S}$ 编码了在给定频率和面内 Bloch 波矢下，整个周期单胞如何把入射通道重排为出射通道。

若把 \cref{eq:rcwa_smatrix} 写成块矩阵形式，
\[
\mathbf{S}=
\begin{bmatrix}
\mathbf{S}_{11} & \mathbf{S}_{12}\\
\mathbf{S}_{21} & \mathbf{S}_{22}
\end{bmatrix},
\]
则 $\mathbf{S}_{11}$ 表示“从上方入射再向上方反射”的块，$\mathbf{S}_{21}$ 表示“从上方入射后向下方透射”的块；其余两块则对应从下方入射的情形。把这层物理含义紧贴在 \cref{eq:rcwa_smatrix} 后面，是为了让后面的 $S_{11}$、$S_{21}$ 不再像纯软件符号。

若只保留零级向上传播与向下传播通道，那么最常见的两个标量散射参量就是
\begin{equation}
S_{11}=\frac{b_{\mathrm{up},0}}{a_{\mathrm{up},0}},
\qquad
S_{21}=\frac{b_{\mathrm{down},0}}{a_{\mathrm{up},0}},
\label{eq:s11_s21_def}
\end{equation}
其中下标 $0$ 表示零级衍射通道，通常要求端口振幅已经按单位功率通量归一化。这样一来，$|S_{11}|^2$ 与 $|S_{21}|^2$ 才能直接对应反射率与透射率的主通道贡献。对存在多个衍射级次的周期结构，更稳妥的做法是分别统计各传播级次的功率通量，再把它们与同一个入射通量相除。把这一步讲清楚的价值在于：后面读反射谱、透射谱、上下出光比例时，读者知道自己读到的是“哪个通道的比值”，而不只是一个黑箱软件吐出的曲线。

对 PCSEL 来说，\cref{eq:rcwa_smatrix} 的威力在于它直接把“内部结构”和“外部辐射通道”连在一起了。于是，上下辐射不对称、偏振选择、法向出光窗口和反射谱中的共振尖峰，都能在同一个框架下被读取。

\begin{RigorousBox}
RCWA 求解的是开放散射问题，而不是封闭本征问题。也正因为如此，它天生适合讨论辐射损耗、反射谱和开放共振。但这并不意味着它自动给出激射结论，因为激射还要求把增益、饱和和注入分布写进系统。开放散射与有源振荡，是相邻但不相同的两个问题层次。
\end{RigorousBox}

\begin{RigorousBox}
RCWA 还有一个经常被忽略、却直接决定收敛质量的技术点：Fourier factorization。对具有介电常数跳变的周期结构，像 $\varepsilon E$ 和 $\varepsilon^{-1}D$ 这样的乘积项并不能一律用“傅里叶系数逐项卷积”来处理。Li 的 Fourier factorization 规则要求根据场分量在界面处的连续性，选择对 $\varepsilon$ 还是对 $\varepsilon^{-1}$ 构造 Toeplitz 矩阵\cite{li1996}。对高折射率对比、矩形孔或金属边界尤其如此；否则，即便傅里叶阶数很高，谱位和线宽也可能出现看似收敛、实则偏差系统性的伪稳定结果。
\end{RigorousBox}

\begin{NoteBox}[breakable]{Fourier Factorization 的数学细节}
设一维周期介电函数 $\varepsilon(x)$ 在单胞内存在间断点。将其展开为傅里叶级数
\begin{equation}
\varepsilon(x) = \sum_{m=-\infty}^{\infty} \hat{\varepsilon}_m \ee^{\ii G_m x},
\qquad
G_m = \frac{2\pi m}{a}.
\end{equation}
截断到 $N$ 阶后，构造 $(2N+1)\times(2N+1)$ Toeplitz 矩阵
\begin{equation}
[\varepsilon]_{mn} = \hat{\varepsilon}_{m-n},
\qquad
m,n = -N,\ldots,N.
\end{equation}
关键问题出现在乘积 $\varepsilon(x)E(x)$ 的傅里叶变换。若电场切向分量 $E_t$ 在界面连续（Maxwell 边界条件），则可用 Laurent 规则
\begin{equation}
\widehat{(\varepsilon E_t)}_m \approx \sum_{n=-N}^{N} [\varepsilon]_{mn} \hat{E}_{t,n}.
\label{eq:fourier_factor_continuous_E}
\end{equation}
但若电位移法向分量 $D_n = \varepsilon E_n$ 在界面连续，而 $E_n$ 本身跳跃，则应对逆介电函数构造矩阵
\begin{equation}
[\varepsilon^{-1}]_{mn} = \widehat{(1/\varepsilon)}_{m-n},
\end{equation}
并使用 inverse rule
\begin{equation}
\hat{E}_{n,m} \approx \sum_{k=-N}^{N} [\varepsilon^{-1}]_{mk} \hat{D}_{n,k},
\qquad
\text{或等价地}\quad
\widehat{(\varepsilon E_n)}_m \approx \left([\varepsilon^{-1}]^{-1}\right)_{mk} \hat{E}_{n,k}.
\label{eq:fourier_factor_discontinuous_E}
\end{equation}
物理上，\cref{eq:fourier_factor_continuous_E} 适用于 TE 偏振的电场切向分量，而 \cref{eq:fourier_factor_discontinuous_E} 适用于 TM 偏振导致的电场法向分量跳跃。数值上，使用错误规则会导致 Gibbs 震荡增强，收敛速度从指数级降为代数级，甚至产生系统性频移。

\begin{NumericalInsightBox}{Factorization 规则的实际影响}
考虑 GaAs/空气二元光栅，折射率对比约 3.4:1。若使用错误的卷积规则：
\begin{itemize}
  \item TM 偏振的共振峰位可能有 5--10 nm 的系统性偏移
  \item 即使傅里叶阶数从 $N=5$ 增加到 $N=21$，偏移仍存在（伪收敛）
  \item 正确的 factorization 可使收敛速度从代数级提升到指数级
\end{itemize}
这就是为什么 commercial RCWA 软件都把 factorization 当成核心技术。对 PCSEL 设计而言，这意味着：如果你发现 RCWA 预测的共振波长与实验测量总有几纳米的固定偏差，在排除工艺误差后，第一个应该检查的就是 Fourier factorization 规则是否正确应用。
\end{NumericalInsightBox}
\end{NoteBox}

把这条规则写成最小可执行形式，就是先记介电函数的 Toeplitz 卷积矩阵为 $[\varepsilon]$，并把基于逆介电函数的矩阵记为 $[\varepsilon^{-1}]$。若某个乘积中的场分量在界面处连续，则可直接使用 Laurent 规则
\begin{equation}
\widehat{\varepsilon E_t}
\approx
[\varepsilon]\,\hat E_t;
\label{eq:rcwa_laurent_rule}
\end{equation}
若法向位移矢量更自然连续，则应使用 inverse rule
\begin{equation}
\widehat{D_n}
=
\widehat{\varepsilon E_n}
\approx
[\varepsilon^{-1}]^{-1}\hat E_n.
\label{eq:rcwa_inverse_rule}
\end{equation}
\cref{eq:rcwa_laurent_rule,eq:rcwa_inverse_rule} 的意义，不是把全部 RCWA 代码细节写完，而是让读者清楚知道：Fourier factorization 不是“截断阶数足够大就自然消失”的数值细节，而是由界面连续性决定的离散化规则。用错卷积矩阵时，结果可能沿着错误的谱位和线宽稳定收敛。

\section{被动谱、共振与极点：谱线为什么会尖起来}
很多读者第一次用 RCWA 时，最直观的输出是一条反射谱或透射谱。看到尖峰或窄谷之后，很容易把它直接称为“高 $Q$ 模”或“低阈值候选模”。这种说法并非总错，但如果不说明其背后的极点结构，就会显得很经验化。

从散射矩阵观点看，开放共振模对应于复频率平面上散射矩阵的极点。通俗地说，如果系统内部存在一个能量停留时间较长的开放模，那么当实频扫描靠近该模的共振位置时，谱线就会出现尖锐变化。谱峰、谱谷、相位跃迁、Fano 线形，本质上都与极点及其和背景散射的干涉有关。

因此，RCWA 中看到尖谱线时，真正应该问的不是“峰高不高”，而是：
\begin{itemize}
  \item 该特征是否随傅里叶阶数和分层数收敛；
  \item 它对应的是哪一支候选模，和 PWEM/CWT 的哪一支模族能对上；
  \item 其线宽对应的被动损耗量级是否与其他方法一致；
  \item 它主要通过哪个辐射通道向外耦合。
\end{itemize}

只有把这些问题答清楚，被动谱才真正转化为器件设计信息。

\section{RCWA 怎样进入光学阈值分析}
前面 \cref{ch:threshold} 已经明确区分了阈值增益和阈值电流。RCWA 在这里的作用，不是直接给出 $I_{\mathrm{th}}$，而是估计周期开放结构层面上的光学阈值增益。最常见的做法，是在有源层引入小信号增益对应的等效虚部，例如写成
\begin{equation}
\varepsilon(\omega)
=
\varepsilon'(\omega)+\ii\varepsilon''_g(g),
\label{eq:rcwa_gain_continuation}
\end{equation}
然后扫描增益参数 $g$，寻找使目标共振从净衰减变为恰好补偿损耗的条件。

如果只停在这一步，读者很容易把“扫到线宽为零”误听成“已经建立了完整激光阈值理论”。教材写法必须把这件事分层说清。对 PCSEL 而言，至少应区分以下四个层次：
\begin{enumerate}
  \item \textbf{冷腔开放共振层}：先在无增益条件下找到复频率极点 $\tilde{\omega}_0=\omega_0-\ii\gamma_0$，它描述的是开放系统如何储能与泄能。
  \item \textbf{线性小信号阈值代理层}：在给定小信号增益模型后继续追踪极点 $\tilde{\omega}(g)$，并用 $\Im \tilde{\omega}(g_{\mathrm{th}})=0$ 估计光学阈值增益。
  \item \textbf{阈值处自洽有源场层}：真正的阈值问题要求场、极化和反转分布同时满足自洽条件，尤其当增益色散明显时，实部与虚部不能被拆成两个彼此独立的旋钮。
  \item \textbf{阈上稳态激光层}：一旦超过阈值，增益饱和、空间烧孔、模式竞争和注入非均匀性都会进入主导地位，这时需要更完整的有源激光形式化，例如基于 Maxwell--Bloch 方程的稳态或时域求解。
\end{enumerate}

本章 RCWA 所做的，严格地说位于第二层。它非常有用，但它仍是\emph{线性小信号阈值代理}，不是阈上稳态激光解。

在这个过程中，可以从几种等价角度理解“达到阈值”：
\begin{itemize}
  \item 被动极点在复频率平面上向实轴靠近；
  \item 某个共振的净线宽收缩到零；
  \item 在小信号线性框架中，目标模的总损耗被增益恰好补偿。
\end{itemize}

这一步得到的是 $g_{\mathrm{th}}$ 这类光学量。它对设计极其有用，因为它告诉我们结构本身对目标模和竞争模提出了多高的增益要求。但必须再次强调，$g_{\mathrm{th}}$ 不是 $I_{\mathrm{th}}$。从前者到后者，还需要材料增益模型、漂移-扩散输运、复合和热反馈。

\paragraph{为什么这只是线性小信号阈值代理}
式 \eqref{eq:rcwa_gain_continuation} 默认我们已经把有源区压缩成了一个“可调复介电常数”模型。只要这一模型仍停留在线性响应层面，那么我们实际上做的是：给开放冷腔模外加一个\emph{尚未饱和}的增益补偿，并观察净衰减是否被抵消。这个判据与 \cref{ch:threshold} 的阈值条件相容，但它并没有告诉我们阈上光场会稳定在什么幅度、是否发生多模竞争、反转是否仍近似均匀，也没有告诉我们空间烧孔会不会把模式排序翻转。

\paragraph{增益色散为什么要求更谨慎}
在许多工程扫描中，人们把 $\varepsilon''_g$ 当成独立参数扫动，而把 $\varepsilon'$ 固定在冷腔值附近。这样做的目的是把问题简化成“先补偿损耗，再看哪一支模最早到零线宽”。作为前端筛选，这种近似经常足够。但若增益谱较窄、工作点离谱峰不远、或载流子引起的色散不可忽略，那么 $\varepsilon'(\omega)$ 与 $\varepsilon''(\omega)$ 由因果材料模型联系在一起，不能任意独立指定。换言之，真正的有源介质不是简单的“只加负虚部”，而是会同时改变频率位置、线宽和模场重叠的复响应系统。

\paragraph{阈值频率并不必然等于冷腔共振频率}
把阈值理解成“在冷腔频率上把线宽调到零”也过于粗糙。更准确的说法是：随着增益参数变化，极点本身会沿复频率平面移动，
\begin{equation}
\tilde{\omega}(g)=\omega_r(g)-\ii\gamma(g).
\end{equation}
阈值条件是 $\gamma(g_{\mathrm{th}})=0$，而真正的阈值频率应写成
\begin{equation}
\omega_{\mathrm{th}}=\Re \tilde{\omega}(g_{\mathrm{th}}),
\end{equation}
它不一定与冷腔共振频率 $\omega_0$ 完全相同。造成这一点的原因包括：增益色散引起的频率拉动、场分布随增益微调而改变，以及不同辐射通道和材料色散共同作用下的复频率漂移。对 PCSEL 来说，这一点尤其重要，因为 $\Gamma$ 点附近的模式排序本来就可能很近，微小的频移就足以改变“哪一支模先达到阈值”的判断。

\paragraph{它与更完整 active lasing formulations 的关系}
从更宽的理论坐标系看，本章 RCWA 增益扫描与更完整的有源激光稳态理论之间是“前端代理”和“自洽主问题”的关系。前者把目标模的阈值判定简化成线性极点继续问题，因此很适合做候选模筛选和参数排序。后者则把 Maxwell 方程与极化、反转和饱和非线性联立起来，直接求解阈值处或阈上稳态激光解；典型代表包括 Maxwell--Bloch 方程的直接时域求解，以及稳态自洽形式如 steady-state \emph{ab initio} laser theory（SALT）\cite{tureci2007salt,tureci2009salt}。本书在这里不展开 SALT 的完整推导，但读者必须知道：RCWA 扫 $\varepsilon''<0$ 到零线宽，\emph{不是}激光理论本体，而是为真正的有源激光问题提供一个高效、可解释、可筛选的线性近似入口。

\section{为什么 RCWA 特别适合做“单胞层面的设计判断”}
一旦接受“RCWA 主要给光学阈值增益和辐射通道信息”，它在工作流中的位置就非常清晰了。它最适合做以下几类判断：

第一，改变孔半径、占空比、刻蚀深度或双晶格扰动时，目标模与竞争模的被动损耗顺序怎样变化。

第二，上下辐射比例、不同偏振通道强弱和法向出光窗口如何随单胞几何变化。

第三，某种外延层微调是否主要通过改变纵向重叠影响辐射，而不是简单改频率。

第四，在后续有限器件仿真之前，哪些参数区间根本不值得继续深挖，因为开放周期单胞层面就已经显示目标模没有优势。

这正是 RCWA 的工程价值所在：它以远低于全尺寸三维时域仿真的代价，把大量“不可能成功”的单胞方案尽早筛掉。

\begin{ExampleBox}
若两种孔形方案在 PWEM 中都显示存在相近的 Gamma 点候选模，但 RCWA 显示方案 A 的目标偏振通道在法向辐射上更有序、而竞争模辐射损耗更大，那么方案 A 更值得进入有限器件验证。此时即便还不能宣布“器件一定低阈值”，也已经完成了一次高价值的前端筛选。
\end{ExampleBox}

\section{RCWA 不能替你解决什么问题}
RCWA 很强，但它解决的是“周期单胞开放散射”问题，而不是“完整器件工作”问题。这个边界必须说得非常硬。

首先，它默认平面内无限周期。因此，有限孔阵边缘泄漏、横向模包络和大面积器件的边界损耗不在默认描述范围内。若这些效应开始主导，必须转向 \cref{ch:fdtd-fem} 的 FDTD/FEM 或更高阶包络模型。

其次，它通常在小信号线性框架里处理增益。于是，增益饱和、空间烧孔和工作点非线性不会自动出现在结果里。

进一步说，即便把小信号增益继续到“零线宽”，也仍未自动得到阈上稳态解。若问题已经进入强色散、强模竞争或反转明显不再静止的区域，就必须转向更完整的有源激光稳态理论，或至少用它们来校准 RCWA 前端筛选的可信边界。

再次，它不包含真实电流路径。均匀增益分布在 RCWA 中只是一个光学代理模型，而不是电注入器件的默认现实。

最后，它也不解决温升问题。很多结构在室温被动谱上看起来很好，一旦工作后温升导致折射率和增益谱漂移，目标模顺序就可能改变。这类问题必须交给热-电-光耦合章节处理。

\begin{PitfallBox}
把 RCWA 单胞反射谱里的最窄特征直接当作“最终器件最低阈值模式”，或把“扫 $\varepsilon''<0$ 直到线宽为零”直接当作“激光阈值理论已经完成”，都是典型的层级偷换。前者忽略了有限尺寸、注入非均匀和温升；后者忽略了增益色散、频率拉动、饱和和模式竞争。
\end{PitfallBox}

\section{收敛与数值卫生：哪些检查是必须交代的}
RCWA 很依赖数值细节。一个结果能不能被相信，不取决于图长得多漂亮，而取决于它是否通过了基本的数值卫生检查。至少应当检查四件事。

第一，傅里叶阶数收敛。目标模的谱位、线宽、反射峰形和上下辐射分配，不能随着阶数增加继续大幅变化。

第二，纵向分层收敛。对刻蚀轮廓、锥形孔或多层外延而言，过粗的分层会把真实边界模糊掉。

第三，能量守恒或功率平衡检查。在无源条件下，反射、透射和吸收的总和应当满足相应守恒关系。

第四，结论对微小参数扰动的鲁棒性。如果某个“最好方案”只在一个极细的参数点上存在，而周围稍微变化就消失，那么它通常不是可靠设计。

这些检查之所以重要，不只是为了数值严谨，也是为了避免把求解器伪影误读成器件物理。

\section{RCWA 在整条工作流中的准确位置}
现在可以把 RCWA 放回全书主线。\cref{ch:pwe} 的 PWEM 给出候选模目录；本章的 RCWA 则对这些候选模做开放周期结构层面的精修。接下来 \cref{ch:fdtd-fem} 的 FDTD/FEM 会把有限尺寸和更真实几何带回来，再由 \cref{ch:epitaxial-optics,ch:qw-gain,ch:electrical,ch:eto-coupling} 把这些光学结果接进外延、增益、电注入和热耦合模型。

也就是说，RCWA 不是终点，而是一个非常关键的中间关口。它回答的问题是：如果先只讨论单胞级开放光学，哪些方案值得继续投入更昂贵的器件级验证。把这个关口守好，后面的工作量会大幅下降；把它省掉，后面的全器件仿真往往会变成盲目试错。

\section*{本章小结}
RCWA 的核心不是某个软件按钮，而是散射矩阵视角：把周期分层开放结构看成一个输入输出系统，直接分析内部模式如何与外界辐射通道交换能量。对 PCSEL 而言，它特别适合处理法向辐射、偏振通道、被动共振和光学阈值增益估计。但本章给出的阈值分析属于\emph{线性小信号阈值代理}：它通过继续追踪开放极点来估计哪一支模最先补偿净损耗，而不是直接求出阈上稳态激光解。只有把这一层级和更完整的有源激光理论分清，RCWA 才能被放在正确位置上使用：作为“开放单胞精修器”和“后续器件仿真的前置筛选器”，而不是最终器件结论的替代品。

\section*{练习题}
\begin{enumerate}
  \item \basicq 用自己的话解释为什么 PWEM 之后自然需要 RCWA，而不是直接跳到器件级电热模型。

  \item \basicq 说明散射矩阵方法为什么在数值上通常比简单传输矩阵更稳定。

  \item \midq 讨论从 RCWA 得到光学阈值增益后，为什么仍然不能直接报告阈值电流。

  \item \midq 解释为什么“在固定 $\varepsilon'$ 下扫描负虚部直到零线宽”属于线性小信号阈值代理，而不是完整激光阈值理论。

  \item \hardq 设计一套你认为最基本的 RCWA 收敛检查清单，并说明每一项各自防止什么误判。
\end{enumerate}

\section*{建议继续阅读}
\begin{itemize}
  \item 下一章将把视角从“周期单胞开放结构”推进到“有限尺寸真实器件”，也就是 FDTD 和 FEM 的领域。读完下一章后再回看本章，会更清楚两者的分工边界。 这条阅读的重点是把本章的输出顺着主线接到后续模块，而不是停成孤立知识点。

  \item 若想打牢开放周期结构的基础，可先对照 \cite{joannopoulos2008,sakoda2005} 中关于周期散射和开放通道的章节阅读。 这条阅读的重点是补足本章关于RCWA、开放单胞谱线与线性阈值代理 的标准背景、经典语言和代表性文献接口。

  \item 若想补充 RCWA 的经典原始文献，可参考 \cite{moharam1981,li1996}。 这条阅读的重点是补足本章关于RCWA、开放单胞谱线与线性阈值代理 的标准背景、经典语言和代表性文献接口。

  \item 若想理解本章线性增益扫描与更完整有源激光理论之间的关系，可进一步阅读 \cite{tureci2007salt,tureci2009salt}。 这条阅读的重点是补足本章关于RCWA、开放单胞谱线与线性阈值代理 的标准背景、经典语言和代表性文献接口。
\end{itemize}


